How can simulations from the posterior and the predictive distributions be im-
plemented? Provide examples of potentially interesting functionals that can be
computed via simulations.

Recall the characterization of the posterior predictive distribution for the new sample \( X_{n+1} \) given the information 
\( X_1 = x_1, \ldots, X_n = x_n \):
\[
f_{X_{n+1} \mid X_1, \ldots, X_n}(x_{n+1} \mid x_1, \ldots, x_n) = \int f_{X_{n+1} \mid \theta}(x_{n+1} \mid \vartheta) \cdot f_{\theta \mid X_1, \ldots, X_n}(\vartheta \mid x_1, \ldots, x_n) \, d\vartheta
\]

Problem: In many modeling situations, we may be able to sample from \( f_{X_i \mid \theta} \) and \( f_{\theta \mid X_1, \ldots, X_n} \), 
but \( f_{X_{n+1} \mid X_1, \ldots, X_n} \) might be too complicated to sample from it directly.

Idea: Indirect sampling using the following Monte Carlo procedure: For each \( r \in \{ 1, \ldots, s \} \):
Sample \( \theta^{(r)} \sim f_{\theta \mid X_1, \ldots, X_n} \), yielding the sample value \( \theta^{(r)} = \vartheta^{(r)} \).
Sample \( X_{n+1}^{(r)} \sim f_{X_{n+1} \mid \theta = \vartheta^{(r)}} \), yielding the sample value \( X_{n+1}^{(r)} = x_{n+1}^{(r)} \).

The sequence \(\{ (\vartheta^{(1)}, x_{n+1}^{(1)}), \ldots, (\vartheta^{(s)}, x_{n+1}^{(s)}) \} \) consists of \( s \) independent samples from 
the joint posterior distribution of \( (\theta, X_{n+1}) \).

The sequence \(\{ x_{n+1}^{(1)}, \ldots, x_{n+1}^{(s)} \} \) consists of \( s \) independent samples from the marginal
posterior distribution of \( X_{n+1} \), i.e., from the \emph{posterior predictive distribution}.

Let \(g\) be a measurable function (satisfying suitable integrability conditions) and let \( \theta^{(1)}, \ldots, \theta^{(s)} \) be i.i.d. samples
from \(f_{\theta \mid X_1 = x_1, \ldots, X_n = x_n}.\) From the law of large numbers (LLN) we obtain:
\(\frac{1}{s} \sum_{r=1}^{s} g(\theta^{(r)}) \xrightarrow{s \to \infty} \mathbb{E}[g(\theta) \mid X_1 = x_1, \ldots, X_n = x_n].\)

Examples of approximations for \( s \to \infty \):
\(\bar{\theta} := \frac{1}{s}\sum_{r=1}^{s} \theta^{(r)} \to \mathbb{E}[\theta \mid X_1, \ldots, X_n]\)


Specify the normal sampling model. Derive a sufficient statistic.

The two-parameter model described by the conditional probability density function
\[
f_{X \mid \theta, \sigma^2}(x \mid \vartheta, \varsigma^2) = \frac{1}{\sqrt{2 \pi \varsigma^2}} \cdot \exp \left( -\frac{1}{2} \left( \frac{y - \vartheta}{\varsigma} \right)^2 \right),
\]

i.e., \(\{ X \mid \theta = \vartheta, \sigma^2 = \varsigma^2 \} \sim \mathcal{N}(\vartheta, \varsigma^2)\) is called the normal model.

Suppose our model is 
\( \{X_1, \ldots, X_n \mid \theta = \vartheta, \sigma^2 = \varsigma^2 \} \) 
i.i.d. normally distributed with parameters \((\vartheta, \varsigma^2)\).

The joint sampling density is given by
\[
f_{x_1, \ldots, x_n \mid \theta, \sigma^2}(x_1, \ldots, x_n \mid \vartheta, \varsigma^2) = \prod_{i=1}^n f_{x_i \mid \theta, \sigma^2}(x_i \mid \vartheta, \varsigma^2) = (2\pi \varsigma^2)^{-n/2} \cdot \exp \left( -\frac{1}{2} \sum_{i=1}^n \left( \frac{x_i - \vartheta}{\varsigma} \right)^2 \right)
\]
\[
= (2 \pi \varsigma^2)^{-n/2} \cdot \exp \left( -\frac{1}{2} \left[ \frac{1}{\varsigma^2} \sum_{i=1}^n x_i^2 - 2 \frac{\vartheta}{\varsigma^2} \sum_{i=1}^n x_i + n \frac{\vartheta^2}{\varsigma^2} \right] \right)
\]
\((\sum_{i=1}^n x_i^2, \sum_{i=1}^n x_i)\) constitutes a two-dimensional sufficient statistic.

Knowing the values of these two quantities is equivalent to knowing the values of
\[
\bar{x} = \frac{1}{n} \sum_{i=1}^n x_i, \quad \text{and} \quad s^2 = \frac{1}{n-1} \sum_{i=1}^n (x_i - \bar{x})^2
\]
\((\bar{x}, s^2)\) also constitutes a sufficient statistic.


Motivate rigorously that the normal is a good candidate for a prior for the mean,
conditionally on the sampling variance.

If 
\(
\{ \theta \mid \sigma^2 = \varsigma^2 \} \sim \mathcal{N}(\mu_0, \tau_0^2), \quad \mu_0 \in \mathbb{R}, \tau_0 > 0
\),
\(
\{ X_1, \ldots, X_n \mid \theta = \vartheta, \sigma^2 = \varsigma^2 \} \sim \mathcal{N}(\vartheta, \varsigma^2),
\)
then
\(
\theta \mid \{X_1 = x_1, \ldots, X_n = x_n, \sigma^2 = \varsigma^2\} \sim \mathcal{N}(\mu_n, \tau_n^2),
\)
where
\(
\tau_n^2 = \frac{1}{\frac{1}{\tau_0^2} + \frac{n}{\varsigma^2}} \quad \text{and} \quad \mu_n = \frac{\frac{1}{\tau_0^2} \mu_0 + \frac{n}{\varsigma^2} \bar{x}}{\frac{1}{\tau_0^2} + \frac{n}{\varsigma^2}}.
\)

ADD REMAINING PROOF (LEFT AS EXERCISE)


Define the notion of precision. State the posterior distribution of the mean, conditionally on the sampling variance. 
How are posterior precision and posterior mean computed?

The (conditional) posterior parameters \( \tau_n^2, \mu_n \),
 combine information from the prior knowledge (captured in the parameters  
\(\tau_0, \mu_0\) of the prior distribution) with information from the sample data:
Posterior Variance and Precision: Note that
\[
\frac{1}{\tau_n^2} = \frac{1}{\tau_0^2} + \frac{n}{\zeta^2}.
\]

The posterior inverse variance is obtained as the sum of prior inverse variance and the inverse of the data variance.
Inverse variance is often called precision. Let:
\(\tilde{\zeta}^2 = 1/\zeta^2\) \sampling precision, i.e., a measure of the distance of the \( x_i \)'s to \( \vartheta \),
\(\tilde{\tau}_0^2 = 1/\tau_0^2\) prior precision,
\(\tilde{\tau}_n^2 = 1/\tau_n^2\) posterior precision
and it holds:
\( \tilde{\tau}_n^2 = \tilde{\tau}_0^2 + \tilde{\zeta}^2 \)

Posterior Mean: Note that
\[
\mu_n = \frac{\tilde{\tau}_0^2}{\tilde{\tau}_0^2 + n \tilde{\zeta}^2} \cdot \mu_0 + \frac{n \tilde{\zeta}^2}{\tilde{\tau}_0^2 + n \tilde{\zeta}^2} \cdot \bar{x}.
\]
The posterior mean is a weighted average of the prior mean and the sample mean.
The weights are 
on the prior mean: \(1/\tau_0^2\), the prior precision
on the sample mean: \( n/\zeta^2\) i.e., the sampling precision of the sample mean.
Assume that the prior mean was based on \(\kappa_0 \) prior observations from the same or similar population \(X_1, \ldots, X_n\).
If we set the prior variance \(\tau_0^2 = \zeta^2/\kappa_0 \) equal to the variance of the mean of prior observations, the posterior mean could be written as
\(
\mu_n = \frac{\kappa_0}{\kappa_0 + n} \mu_0 + \frac{n}{\kappa_0 + n} \bar{x}.
\).


 What is the predictive distribution in the normal model?

Note that we can assume \( \{ X_{n+1} | \theta = \vartheta, \sigma^2 = \varsigma^2 \} \sim \mathcal{N}(\vartheta, \varsigma^2) \)
\( \iff X_{n+1} = \theta + \epsilon, \; \{ \epsilon | \theta = \vartheta, \sigma^2 = \varsigma^2 \} \sim \mathcal{N}(0, \varsigma^2) \),
and \( \theta | \sigma^2 \text{ and } \epsilon | \sigma^2 \) are conditionally independent.

Then wen can compute posterior mean and variance of \( X_{n+1} \): 

\(
\mathbb{E}[X_{n+1}|X_1 = x_1, \ldots, X_n = x_n, \sigma^2 = \varsigma^2] = \mathbb{E}[\theta + \epsilon|X_1 = x_1, \ldots, X_n = x_n, \sigma^2 = \varsigma^2]
\)
\(
= \mathbb{E}[\theta|X_1 = x_1, \ldots, X_n = x_n, \sigma^2 = \varsigma^2] + \mathbb{E}[\epsilon|X_1 = x_1, \ldots, X_n = x_n, \sigma^2 = \varsigma^2]
\)
\(
= \mu_n + 0 = \mu_n
\)
\(
\operatorname{Var}[X_{n+1}|X_1 = x_1, \ldots, X_n = x_n, \sigma^2 = \varsigma^2] = \operatorname{Var}[\theta + \epsilon|X_1 = x_1, \ldots, X_n = x_n, \sigma^2 = \varsigma^2]
\)
\(
= \tau_n^2 + \varsigma^2.
\)
It can be shown that the predictive posterior distribution is
\( \{ X_{n+1}|X_1 = x_1, \ldots, X_n = x_n, \sigma^2 = \varsigma^2 \} \sim \text{N}(\mu_n, \tau^2_n + \varepsilon^2 \)

In the limit,
we become increasingly certain about \( \theta \), i.e., the posterior variance \( \tau_n^2 \) approaches zero,
albeit not reducing \( \varepsilon^2 \) (sampling variance); which is therefore a lower bound of our uncertainty.


If we consider the joint inference for mean and variance of the sampling distribu-
tion, what is a convenient prior for the variance?


Consider the following two-parameter sampling model:

\( 1/\sigma^2 \sim \Gamma(v_0/2, v_0\sigma_0^2/2) \)
\( \{\theta\} \sim N(\mu_0, \tau_0^2) \)
\( \{X_1, ..., X_n | \theta = \vartheta, \sigma^2 = \varsigma^2 \} \sim i.i.d. \ N(\theta, \sigma^2) \)

DON'T INCLUDE THIS ITS MORE THAN NECESSARY

Using the independence of priors for \( \theta \) and \( \sigma^2 \), one can show for the precision \( \tilde{\sigma}^2 \):
\[ f_{\sigma^2 | \theta, x_1, ..., x_n}(\varsigma^2 | \theta, x_1, ..., x_n) \propto \left( \varsigma^2 \right)^{n/2} \exp \left\{ -\frac{1}{2}\varsigma^2 \sum_{i=1}^{n} (x_i - \theta)^2 \right\} \cdot \left( \varsigma^2 \right)^{\nu_0/2 - 1} \exp \left\{ -\frac{1}{2}\varsigma^2 \nu_0 \sigma_0^2 \right\} \]
\[ = \left( \varsigma^2 \right)^{\frac{(\nu_0 + n)}{2} - 1} \cdot \exp \left\{ -\frac{1}{2}\varsigma^2 \cdot \left( \nu_0 \sigma_0^2 + \sum_{i=1}^{n} (x_i - \theta)^2 \right) \right\} \]
Identifying a gamma density, we obtain
\[ \{\sigma^2 | \theta = \vartheta, X_1 = x_1, ..., X_n = x_n \} \sim IG(\nu_n/2, \nu_n \sigma_n^2(\theta)/2), \]
where
\[ \nu_n = \nu_0 + n, \quad \sigma_n^2(\theta) = \frac{1}{\nu_n} [\nu_0 \sigma_0^2 + n s_n^2(\theta)], \]
and \( s_n^2(\theta) = \sum_{i=1}^{n} (x_i - \theta)^2 / n \) is the unbiased estimator of \( \sigma^2 \) given \( \theta = \vartheta \).

The benefit is that the posterior distributions match the priors for \( \theta, \sigma^2 \).



Explain in detail how Gibbs sampling can be applied
in the independent case for the Normal model. Extend this to the multivariate case.



Problem:
The marginal distribution of \( \sigma^2 \) given the data is in this case not an inverse-gamma distribution.
It is even not another standard distribution from which we could easily sample!

In the independent prior normal model, we can use the full conditional distributions of the parameters to sample from the 
joint posterior distribution using the Gibbs sampler:

Gibbs Sampling in the Normal Model

Input: Initial values \(\Phi^{(0)} = \{\theta^{(0)}, \sigma^{2(0)}\}\); Set \(s = 1\).

Algorithm:
Step 1: Sample \(\theta^{s} \sim f_{\theta | \sigma^2, x_1,...,x_n}\) (normal)
Step 2: Sample \(\sigma^{2(s)} \sim f_{\sigma^2 | \theta^{(s)}, x_1,...,x_n}\) (gamma)
Set \(s = s + 1\), return to Step 1.
Output: Dependent sequence of parameter values
\[
\Phi^{(1)} = \{\theta^{(1)}, \sigma^{2(1)}\}, \Phi^{(2)} = \{\theta^{(2)}, \sigma^{2(2)}\}, ...
\]

Multivariate:
If \(X_1,...,X_n\) are independent samples from a multivariate normal distribution with mean \(\theta\), then a convenient prior distribution for 
\(\theta\) is a multivariate normal distribution \(\theta \sim N(\mu_0, \Lambda_0)\).

In this case, one can show that the full conditional distribution is given by:
\[
(\theta | X_1,...,X_n, \Sigma = \varsigma) \sim N(\mu_n, \Lambda_n),
\]
where \(\Lambda_n = (\Lambda_0 + n\varsigma^{-1})^{-1}\) and \(\mu_n = \Lambda_n (\Lambda_0 \mu_0 + n\varsigma^{-1}\bar{x})\).

A multivariate generalization of the gamma distribution is the so-called generalized Wishart distribution. 
Analogous to the univariate case, a generalized inverse-Wishart distribution prior for the covariance matrix results in a 
generalized inverse-Wishart full conditional distribution.

\iffalse
Consider the following two-parameter sampling model:

\( 1/\sigma^2 \sim \Gamma(v_0/2, v_0\sigma_0^2/2) \)
\( \{\theta | \sigma^2 = s^2 \} \sim N(\mu_0, s^2/\kappa_0) \)
\( \{X_1, ..., X_n | \theta = \vartheta, \sigma^2 = \varsigma^2 \} \sim i.i.d. \ N(\theta, \sigma^2) \)

The joint posterior density of \( \theta \) and \( \sigma^2 \) can be decomposed as
\[ f_{\theta, \sigma^2 | x_1, ..., x_n} = f_{\theta | \sigma^2, x_1, ..., x_n} \cdot f_{\sigma^2 | x_1, ..., x_n} \]

From a previous section, we know that
\( \{\theta | X_1 = x_1, ..., X_n = x_n, \sigma^2 = s^2 \} \sim N(\mu_n, s^2/\kappa_n) \)
where
\[ \kappa_n = \kappa_0 + n \]
and
\[ \mu_n = \frac{\kappa_0\mu_0 + n\bar{x}}{\kappa_n} \]
Based on
\[ f_{\sigma^2 | x_1, ..., x_n}(s^2 | x_1, ..., x_n) \propto f_{\sigma^2}(s^2) \cdot \int_{-\infty}^{\infty} f_{x_1, ..., x_n | \theta, \sigma^2}(x_1, ..., x_n | \vartheta, s^2) \cdot f_{\theta | \sigma^2}(\vartheta | s^2) d\vartheta, \]
one may show in tedious computations that
\[ \{1/\sigma^2 | X_1 = x_1, ..., X_n = x_n \} \sim \Gamma(\nu_n/2, \nu_n \sigma_n^2/2), \]
where
\[ \nu_n = \nu_0 + n \]
and
\[ \sigma_n^2 = \frac{\nu_0 \sigma_0^2 + (n-1)s^2 + \frac{\kappa_0 n}{\kappa_0 + n}(\bar{x} - \mu_0)^2}{\nu_n}. \]
\fi
